
les
flèches écrites ici sont des homomorphismes de restriction, et on utilise le fait que @@M0@@ est un isomorphisme d'après le théorème de changement de base propre (SGA 4 XII 5.5). Je dis que \emph{le carré} (7.13.2) \emph{est commutatif}. La démonstration est essentiellement la même que celle qui établit la commutativité de (7.10.2); il suffit ici d'utiliser la compatibilité de la formation des classes de Chern avec les images inverses, dans le cas des morphismes @@M1@@ et @@M2@@.

\textbf{7.13.3.} Notons que le carré (7.13.2) provient par passage à la limite projective des carrés analogues, où la cohomologie à coefficients dans le faisceau @@M3@@-adique @@M4@@ est remplacée par celle à coefficients dans les faisceaux de torsion @@M5@@ (où @@M6@@ parcourt les puissances de @@M7@@). Donc la commutativité de (7.13.2) ne fait qu'exprimer la commutativité des carrés faisant intervenir les classes de Chern mod. @@M8@@, pour toute puissance @@M9@@ de @@M10@@. Choisissons une clôture algébrique @@M11@@ de @@M12@@, d'où une clôture algébrique @@M13@@ de @@M14@@ (savoir le corps résiduel de la clôture normale de @@M15@@ dans @@M16@@). On peut, pour toute sous-extension finie @@M17@@ de @@M18@@, considérer le diagramme commutatif analogue (pour @@M19@@ fixé) relatif à la situation déduite de @@M20@@ par le changement de base @@M21@@, @@M22@@ étant le normalisé de @@M23@@ dans @@M24@@. On notera pour ceci que, par Krull--Akizuki, et @@M25@@ étant hensélien, @@M26@@ est bien un trait; de plus, il faudra supposer ici que les fibres de @@M27@@ sur @@M28@@ sont géométriquement régulières, i.e. @@M29@@ \emph{lisse} sur @@M30@@, de sorte que cette hypothèse se conserve par changement de base. On trouve ainsi, pour @@M31@@ variable, un système inductif de carrés commutatifs. Passant à la limite inductive et utilisant (SGA 4 VII 5.8) et (IV 3.2), on trouve un carré commutatif (cf. 7.16) :
